%! root=./main.tex
\section{Evaluation}\label{sec:eval}

\subsection{Efficiency}
Here we present the efficiency evaluation result.

%! root=../main.tex
\begin{table}[]
\caption{Runtime overhead on SPEC CPU 2017 benchmarks.}~\label{tab:spec2017}
    \resizebox{.8\linewidth}{!}{
        \begin{tabular}{@{}llll@{}}
        \toprule
            Benchmark & Overhead & Benchmark & Overhead \\ \midrule
            500.perlbench\_r & 31.7204\% & 525.x264\_r & 7.7399\% \\
            502.gcc\_r & 23.1884\% & 531.deepsjeng\_r & 26.5152\% \\
            505.mcf\_r & 15.1460\% & 541.leela\_r & 34.1421\% \\
            520.omnetpp\_r & 28.5714\% & 548.exchange2\_r & 2.0045\% \\
            523.xalancbmk\_r & 19.7260\% & 557.xz\_r & 1.6529\% \\ \bottomrule
        \end{tabular}
    }
\end{table}

\subsection{Security}
Here we present the security evaluation and policy writing instructions. Firstly, we went over to confirm


\subsection{Case studies}
Here we present/discussion subtle/easy-to-overlook bugs (or design flaws) might have overlooked otherwise.

\kjee{tentative items; we hope to have more here.}

\heading{Address translation.}
While describing the $i'2i$ function used in the proof, we briefly mentioned an
unintuitive behavior that causes jumping to the old text segment to be permitted since old text segment indices are left unchanged,
but gave an explanation on why this is still fine.
However, leaving the analogous behavior present in the $i2i'$ function used in the rewriter implementation would invalidate the proof.
Instead, when any index in the new text segment is provided as input to $i2i'$, we must translate it to $ai$.
To show why this is necessary, consider a policy that permits an instruction to jump to an index $i$ outside of the old text segment
that happens to be part of the new text segment after rewriting.
A hash table for this instruction would contain $i2i'(i)$ as one of its entries, which equals $i$ under the naive definition,
as $i$ is not in the old text segment.
Since $i$ could be any index of the new text segment, the destination does not have to be the start of another block.
If the end of a dynamic control flow block was at $i$, a trace could execute a dynamic branch without executing the guard code
that constrains the instruction to permitted destinations. Thus, $i2i'$ cannot leave new text segment indices alone
and instead must choose a different policy permitted destination to place in the table, such as $ai$.

\heading{Branching to PC.}
Most instructions other than those in the data processing or load/store categories are defined to be \unpredictable{} 
when using PC as a register operand, as there is no practical reason to use PC in these instructions.
It is reasonable to assume that \texttt{BX pc} (branch to the program counter) would be included, but it is actually
a valid instruction used in Thumb mode.
In ARM mode, this instruction behaves exactly the same as a static branch with an offset of 8, so there is no reason to use it
(the GNU assembler will explicitly warn that "use of r15 in bx in ARM mode is not really useful"
if you try).
If the rewriter mistakenly handles this instruction like other \texttt{BX} instructions,
then the corresponding block will start with a \texttt{UBFX} that writes to PC, which is \unpredictable.

\heading{Unaligned PC Load.}
\mxz{todo}
